% --- LaTeX Presentation Template - Complete Example ---

% --- Set document class ---

% Remove "handout" when presenting to include pauses
\documentclass[aspectratio=169, dvipsnames]{beamer}
\usetheme{default}

% Make content that is hidden by pauses "transparent"
\setbeamercovered{transparent}

% --- Slide layout settings ---

% Set line spacing
\renewcommand{\baselinestretch}{1.15}

% Set left and right text margins
\setbeamersize{text margin left=12mm, text margin right=12mm}

% Remove navigation symbols
\setbeamertemplate{navigation symbols}{}

% Allow local line spacing changes
\usepackage{setspace}

% Change itemized list bullets to circles
\setbeamertemplate{itemize item}{$\bullet$}
\setbeamertemplate{itemize subitem}{$\circ$}

% --- Color and font settings ---

\usepackage{xcolor}
\usepackage{graphicx}
\usepackage{subfigure}
\usepackage{booktabs}
\usepackage{multirow}
\usepackage{algorithm}
\usepackage{algpseudocode}
\usepackage{float}
\usepackage{natbib}
\usepackage{amsmath, amsfonts, amssymb}
\usepackage{dsfont}
\usepackage{tikz}
\usepackage{mathtools}
\usepackage{euscript}

% Bibliography style
\bibliographystyle{plain}

% Slide title background color
\setbeamercolor{frametitle}{bg=blue!10, fg=black}

% Colors: block title and background
\setbeamercolor{block title}{fg=white, bg=cyan!50!black}
\setbeamercolor{block body}{bg=cyan!10}

% --- Title and author information ---

\title{Example Beamer Presentation}
\author{Your Name}
\date{November 2025}

% --- Custom footline ---
\setbeamertemplate{footline}{%
    \leavevmode%
    \hbox{%
        \begin{beamercolorbox}[wd=.333333\paperwidth,ht=2.5ex,dp=1ex,center]{author in head/foot}%
            \usebeamerfont{author in head/foot}\insertshortauthor
        \end{beamercolorbox}%
        \begin{beamercolorbox}[wd=.333333\paperwidth,ht=2.5ex,dp=1ex,center]{title in head/foot}%
            \usebeamerfont{title in head/foot}\insertshorttitle
        \end{beamercolorbox}%
        \begin{beamercolorbox}[wd=.333333\paperwidth,ht=2.5ex,dp=1ex,center]{date in head/foot}%
            \insertframenumber{} / \inserttotalframenumber
        \end{beamercolorbox}%
    }%
    \vskip0pt%
}

% --- Document begins here ---

\begin{document}

% Title slide
\begin{frame}
    \titlepage
\end{frame}

% Table of contents
\begin{frame}{Outline}
    \tableofcontents
\end{frame}

% --- Section 1: Introduction ---

\section{Introduction}

\begin{frame}{Welcome}
    \begin{itemize}
        \item This is an example presentation using the Beamer template
        \pause
        \item Demonstrates various features and content types
        \pause
        \item Shows best practices for academic presentations
    \end{itemize}
\end{frame}

\begin{frame}{Key Features}
    \begin{columns}
        \column{0.48\textwidth}
        \textbf{Design Features:}
        \begin{itemize}
            \item Clean 16:9 layout
            \item Custom footline
            \item Progressive reveal
            \item Color themes
        \end{itemize}
        
        \column{0.48\textwidth}
        \textbf{Content Support:}
        \begin{itemize}
            \item Mathematics
            \item Algorithms
            \item Figures \& tables
            \item Citations
        \end{itemize}
    \end{columns}
\end{frame}

% --- Section 2: Mathematical Content ---

\section{Mathematical Content}

\begin{frame}{Equations}
    Inline equations are easy: $E = mc^2$
    
    \vspace{1em}
    
    Display equations with numbering:
    \begin{equation}
        \int_{-\infty}^{\infty} e^{-x^2} dx = \sqrt{\pi}
    \end{equation}
    
    \vspace{1em}
    
    Multiple equations:
    \begin{align}
        f(x) &= x^2 + 2x + 1 \\
        g(x) &= (x+1)^2
    \end{align}
\end{frame}

\begin{frame}{Advanced Math}
    Matrix notation:
    \begin{equation}
        \mathbf{A} = \begin{bmatrix}
            a_{11} & a_{12} & \cdots & a_{1n} \\
            a_{21} & a_{22} & \cdots & a_{2n} \\
            \vdots & \vdots & \ddots & \vdots \\
            a_{m1} & a_{m2} & \cdots & a_{mn}
        \end{bmatrix}
    \end{equation}
    
    \vspace{1em}
    
    Probability notation:
    \begin{equation}
        P(X = x) = \frac{e^{-\lambda}\lambda^x}{x!}, \quad x \in \{0, 1, 2, \ldots\}
    \end{equation}
\end{frame}

% --- Section 3: Algorithms ---

\section{Algorithms}

\begin{frame}{Algorithm Example}
    \begin{algorithm}[H]
        \caption{Gradient Descent}
        \begin{algorithmic}[1]
            \State \textbf{Input:} Initial point $x_0$, learning rate $\alpha$, iterations $T$
            \State \textbf{Output:} Optimized point $x^*$
            \For{$t = 1$ to $T$}
                \State Compute gradient $\nabla f(x_t)$
                \State Update: $x_{t+1} \gets x_t - \alpha \nabla f(x_t)$
            \EndFor
            \State \textbf{return} $x^* = x_T$
        \end{algorithmic}
    \end{algorithm}
\end{frame}

% --- Section 4: Visual Content ---

\section{Visual Content}

\begin{frame}{Tables}
    \begin{table}
        \centering
        \caption{Experimental Results}
        \begin{tabular}{lccc}
            \toprule
            \textbf{Method} & \textbf{Accuracy} & \textbf{Precision} & \textbf{Recall} \\
            \midrule
            Baseline & 85.3 & 82.1 & 87.2 \\
            Method A & 89.7 & 88.3 & 90.5 \\
            \textbf{Our Method} & \textbf{92.4} & \textbf{91.8} & \textbf{93.1} \\
            \bottomrule
        \end{tabular}
    \end{table}
\end{frame}

\begin{frame}{Block Environments}
    \begin{block}{Definition}
        A \textbf{metric space} is a set $X$ equipped with a function $d: X \times X \to \mathbb{R}$ satisfying certain properties.
    \end{block}
    
    \vspace{1em}
    
    \begin{alertblock}{Important Note}
        Always verify convergence conditions before applying iterative methods.
    \end{alertblock}
    
    \vspace{1em}
    
    \begin{exampleblock}{Example}
        The Euclidean space $\mathbb{R}^n$ with distance $d(\mathbf{x}, \mathbf{y}) = \|\mathbf{x} - \mathbf{y}\|_2$ is a metric space.
    \end{exampleblock}
\end{frame}

% --- Section 5: Conclusion ---

\section{Conclusion}

\begin{frame}{Summary}
    \begin{itemize}
        \item This template provides a clean, professional presentation format
        \item Supports mathematics, algorithms, figures, and tables
        \item Easy to customize colors and layout
        \item Suitable for academic talks and conferences
    \end{itemize}
    
    \vspace{2em}
    
    \begin{center}
        \Large Thank you for your attention!
    \end{center}
\end{frame}

\begin{frame}{Questions?}
    \begin{center}
        \Huge Questions?
        
        \vspace{2em}
        
        \large Contact: your.email@university.edu
    \end{center}
\end{frame}

\end{document}
